\begin{remark*}[Orientation]
Dilation and displacement are affine operations on sets.  Translation preserves
the order of all comparisons, positive dilation preserves order, and negative
dilation reverses order.
\end{remark*}

\begin{definitionbox}{Definition (Displacement, Dilation, and Reflection)}
\begin{definition}[Displacement, Dilation, and Reflection]
\label{def:displacement-dilation-reflection}
Let $A\subseteq\mathbb{R}$. For $c,\lambda\in\mathbb{R}$, the
\textbf{displacement} of $A$ by $c$, the \textbf{dilation} of $A$ by
$\lambda$, and the \textbf{reflection} of $A$ through the origin are
defined, respectively, by
\[
A+c:=\{a+c:a\in A\},
\qquad
\lambda A:=\{\lambda a:a\in A\},
\qquad
-A:=\{-a:a\in A\}.
\]
\end{definition}
\end{definitionbox}

\begin{remark*}[Standard quantified statement]
\[
x\in A+c\Longleftrightarrow(\exists a\in A)(x=a+c),
\qquad
x\in\lambda A\Longleftrightarrow(\exists a\in A)(x=\lambda a),
\qquad
x\in -A\Longleftrightarrow(\exists a\in A)(x=-a).
\]
\end{remark*}

\begin{remark*}[Predicate reading]
\[
\begin{aligned}
&A\subseteq\mathbb{R},\qquad A+c=\{a+c:a\in A\},\qquad
\lambda A=\{\lambda a:a\in A\},\qquad -A=\{-a:a\in A\},\\
&\varphi_c\colon A\to A+c,\qquad
\varphi_\lambda\colon A\to\lambda A,\qquad
\varphi_-\colon A\to -A,\\
&\operatorname{FunctionalRelation}(\varphi_c,A,A+c),\qquad
\operatorname{FunctionalRelation}(\varphi_\lambda,A,\lambda A),\qquad
\operatorname{FunctionalRelation}(\varphi_-,A,-A),\\
&\varphi_c(a)=a+c,\qquad \varphi_\lambda(a)=\lambda a,\qquad \varphi_-(a)=-a,\\
&\operatorname{ReplacementImage}(A+c,A,\varphi_c),\qquad
\operatorname{ReplacementImage}(\lambda A,A,\varphi_\lambda),\qquad
\operatorname{ReplacementImage}(-A,A,\varphi_-).
\end{aligned}
\]
\end{remark*}

\begin{remark*}[Negated quantified statement]
\[
x\notin A+c\Longleftrightarrow(\forall a\in A)(x\neq a+c),
\qquad
x\notin\lambda A\Longleftrightarrow(\forall a\in A)(x\neq\lambda a).
\]
\end{remark*}

\begin{remark*}[Negation predicate reading]
\[
\begin{aligned}
&A\subseteq\mathbb{R},\qquad A+c=\{a+c:a\in A\},\qquad
\varphi_c\colon A\to A+c,\qquad
\operatorname{FunctionalRelation}(\varphi_c,A,A+c),\qquad \varphi_c(a)=a+c,\\
&\neg\operatorname{ReplacementImage}(B,A,\varphi_c)
\end{aligned}
\quad\text{means}\quad
(\exists x)\bigl[x\in B\triangle(A+c)\bigr],
\]
and similarly for dilation and reflection.
\end{remark*}

\begin{remark*}[Failure modes]
\begin{description}
\item[Exposition.]
The definition fails only if the image set omits an actual transformed point
or contains a point that is not produced by transforming an element of $A$.


  \item[\operatorname{FunctionalRelation} / \operatorname{ReplacementImage}.]
  \[
x\notin A+c\Longleftrightarrow(\forall a\in A)(x\neq a+c),
\qquad
x\notin\lambda A\Longleftrightarrow(\forall a\in A)(x\neq\lambda a).
\]
\[
\begin{aligned}
&A\subseteq\mathbb{R},\qquad A+c=\{a+c:a\in A\},\qquad
\varphi_c\colon A\to A+c,\qquad
\operatorname{FunctionalRelation}(\varphi_c,A,A+c),\qquad \varphi_c(a)=a+c,\\
&\neg\operatorname{ReplacementImage}(B,A,\varphi_c)
\end{aligned}
\quad\text{means}\quad
(\exists x)\bigl[x\in B\triangle(A+c)\bigr],
\]
\end{description}
\end{remark*}

\begin{remark*}[Failure modes]
\begin{description}
\item[Exposition.]
Every membership claim in a transformed set decomposes into an original element
and the transformation rule that produced the new point.


  \item[\operatorname{FunctionalRelation} / \operatorname{ReplacementImage}.]
  \[
x\notin A+c\Longleftrightarrow(\forall a\in A)(x\neq a+c),
\qquad
x\notin\lambda A\Longleftrightarrow(\forall a\in A)(x\neq\lambda a).
\]
\[
\begin{aligned}
&A\subseteq\mathbb{R},\qquad A+c=\{a+c:a\in A\},\qquad
\varphi_c\colon A\to A+c,\qquad
\operatorname{FunctionalRelation}(\varphi_c,A,A+c),\qquad \varphi_c(a)=a+c,\\
&\neg\operatorname{ReplacementImage}(B,A,\varphi_c)
\end{aligned}
\quad\text{means}\quad
(\exists x)\bigl[x\in B\triangle(A+c)\bigr],
\]
\end{description}
\end{remark*}

\begin{remark*}[Interpretation]
Displacement slides the set, dilation stretches or collapses it, and reflection
is the special dilation that reverses the real line.
\end{remark*}

\begin{dependencies}
\begin{itemize}
  \item \hyperref[def:reals]{Real Numbers}
  \item \hyperref[def:subset]{Subset}
  \item \hyperref[def:function]{Function}
  \item \hyperref[def:functional-relation]{Functional Relation}
  \item \hyperref[def:image-set]{Image Set}
  \item \hyperref[def:addition-on-r]{Addition on $\mathbb{R}$}
  \item \hyperref[def:subtraction-on-r]{Subtraction on $\mathbb{R}$}
  \item \hyperref[def:multiplication-on-r]{Multiplication on $\mathbb{R}$}
\end{itemize}
\end{dependencies}

\begin{propositionbox}{Proposition (Translation Preserves Upper Bounds)}
\begin{proposition}[Translation Preserves Upper Bounds]
\label{prop:translation-preserves-upper-bounds}
Let $A\subseteq\mathbb{R}$ and let $c,u\in\mathbb{R}$.  If $u$ is an
upper bound of $A$, then $u+c$ is an upper bound of $A+c$; that is,
\[
\left(\forall a\in A,\ a\leq u\right)
\Longrightarrow
\left(\forall x\in A+c,\ x\leq u+c\right).
\]
\hyperref[prf:translation-preserves-upper-bounds]{Go to proof.}
\end{proposition}
\end{propositionbox}

\begin{remark*}[Standard quantified statement]
\[
(\forall a\in A)(a\leq u)
\Longrightarrow
(\forall x\in A+c)(x\leq u+c).
\]
\end{remark*}

\begin{remark*}[Predicate reading]
\[
\begin{aligned}
&P=\mathsf{OrderedSet}(\mathbb{R},\leq),\qquad
A\subseteq\mathbb{R},\qquad A+c=\{a+c:a\in A\},\\
&\operatorname{UpperBound}(u,A,P)\Longrightarrow\operatorname{UpperBound}(u+c,A+c,P).
\end{aligned}
\]
\end{remark*}

\begin{remark*}[Negated quantified statement]
\[
(\forall a\in A)(a\leq u),\qquad
(\exists y\in A+c)(u+c<y).
\]
\end{remark*}

\begin{remark*}[Negation predicate reading]
\[
\begin{aligned}
&P=\mathsf{OrderedSet}(\mathbb{R},\leq),\qquad
A\subseteq\mathbb{R},\qquad A+c=\{a+c:a\in A\},\\
&\operatorname{UpperBound}(u,A,P)\land\neg\operatorname{UpperBound}(u+c,A+c,P).
\end{aligned}
\]
\end{remark*}

\begin{remark*}[Failure modes]
\begin{description}
\item[Exposition.]
A failure would require a translated point above the translated upper bound.


  \item[\operatorname{UpperBound}.]
  \[
(\forall a\in A)(a\leq u),\qquad
(\exists y\in A+c)(u+c<y).
\]
\[
\begin{aligned}
&P=\mathsf{OrderedSet}(\mathbb{R},\leq),\qquad
A\subseteq\mathbb{R},\qquad A+c=\{a+c:a\in A\},\\
&\operatorname{UpperBound}(u,A,P)\land\neg\operatorname{UpperBound}(u+c,A+c,P).
\end{aligned}
\]
\end{description}
\end{remark*}

\begin{remark*}[Failure modes]
\begin{description}
\item[Exposition.]
Write $y=a+c$ with $a\in A$.  Then $a\leq u$ gives $a+c\leq u+c$.


  \item[\operatorname{UpperBound}.]
  \[
(\forall a\in A)(a\leq u),\qquad
(\exists y\in A+c)(u+c<y).
\]
\[
\begin{aligned}
&P=\mathsf{OrderedSet}(\mathbb{R},\leq),\qquad
A\subseteq\mathbb{R},\qquad A+c=\{a+c:a\in A\},\\
&\operatorname{UpperBound}(u,A,P)\land\neg\operatorname{UpperBound}(u+c,A+c,P).
\end{aligned}
\]
\end{description}
\end{remark*}

\begin{remark*}[Interpretation]
Translation shifts the set and its upper bound by the same amount, so the
relative order between points and the bound is unchanged.
\end{remark*}

\begin{dependencies}
\begin{itemize}
  \item \hyperref[def:reals]{Real Numbers}
  \item \hyperref[def:ordered-set]{Ordered Set}
  \item \hyperref[def:subset]{Subset}
  \item \hyperref[def:addition-on-r]{Addition on $\mathbb{R}$}
  \item \hyperref[def:displacement-dilation-reflection]{Displacement, Dilation, and Reflection}
  \item \hyperref[def:real-upper-bound]{Upper Bound}
\end{itemize}
\end{dependencies}

\begin{propositionbox}{Proposition (Translation Preserves Lower Bounds)}
\begin{proposition}[Translation Preserves Lower Bounds]
\label{prop:translation-preserves-lower-bounds}
Let $A\subseteq\mathbb{R}$ and let $c,\ell\in\mathbb{R}$.  If $\ell$ is
a lower bound of $A$, then $\ell+c$ is a lower bound of $A+c$; that is,
\[
\left(\forall a\in A,\ \ell\leq a\right)
\Longrightarrow
\left(\forall x\in A+c,\ \ell+c\leq x\right).
\]
\hyperref[prf:translation-preserves-lower-bounds]{Go to proof.}
\end{proposition}
\end{propositionbox}

\begin{remark*}[Standard quantified statement]
\[
(\forall a\in A)(\ell\leq a)
\Longrightarrow
(\forall x\in A+c)(\ell+c\leq x).
\]
\end{remark*}

\begin{remark*}[Predicate reading]
\[
\begin{aligned}
&P=\mathsf{OrderedSet}(\mathbb{R},\leq),\qquad
A\subseteq\mathbb{R},\qquad A+c=\{a+c:a\in A\},\\
&\operatorname{LowerBound}(\ell,A,P)\Longrightarrow\operatorname{LowerBound}(\ell+c,A+c,P).
\end{aligned}
\]
\end{remark*}

\begin{remark*}[Negated quantified statement]
\[
(\forall a\in A)(\ell\leq a),\qquad
(\exists y\in A+c)(y<\ell+c).
\]
\end{remark*}

\begin{remark*}[Negation predicate reading]
\[
\begin{aligned}
&P=\mathsf{OrderedSet}(\mathbb{R},\leq),\qquad
A\subseteq\mathbb{R},\qquad A+c=\{a+c:a\in A\},\\
&\operatorname{LowerBound}(\ell,A,P)\land\neg\operatorname{LowerBound}(\ell+c,A+c,P).
\end{aligned}
\]
\end{remark*}

\begin{remark*}[Failure modes]
\begin{description}
\item[Exposition.]
A failure would require a translated point below the translated lower bound.


  \item[\operatorname{LowerBound}.]
  \[
(\forall a\in A)(\ell\leq a),\qquad
(\exists y\in A+c)(y<\ell+c).
\]
\[
\begin{aligned}
&P=\mathsf{OrderedSet}(\mathbb{R},\leq),\qquad
A\subseteq\mathbb{R},\qquad A+c=\{a+c:a\in A\},\\
&\operatorname{LowerBound}(\ell,A,P)\land\neg\operatorname{LowerBound}(\ell+c,A+c,P).
\end{aligned}
\]
\end{description}
\end{remark*}

\begin{remark*}[Failure modes]
\begin{description}
\item[Exposition.]
Write $y=a+c$ with $a\in A$.  Then $\ell\leq a$ gives $\ell+c\leq a+c$.


  \item[\operatorname{LowerBound}.]
  \[
(\forall a\in A)(\ell\leq a),\qquad
(\exists y\in A+c)(y<\ell+c).
\]
\[
\begin{aligned}
&P=\mathsf{OrderedSet}(\mathbb{R},\leq),\qquad
A\subseteq\mathbb{R},\qquad A+c=\{a+c:a\in A\},\\
&\operatorname{LowerBound}(\ell,A,P)\land\neg\operatorname{LowerBound}(\ell+c,A+c,P).
\end{aligned}
\]
\end{description}
\end{remark*}

\begin{remark*}[Interpretation]
Translation preserves lower-bound information because adding the same constant
to both sides preserves the order relation.
\end{remark*}

\begin{dependencies}
\begin{itemize}
  \item \hyperref[def:reals]{Real Numbers}
  \item \hyperref[def:ordered-set]{Ordered Set}
  \item \hyperref[def:subset]{Subset}
  \item \hyperref[def:addition-on-r]{Addition on $\mathbb{R}$}
  \item \hyperref[def:displacement-dilation-reflection]{Displacement, Dilation, and Reflection}
  \item \hyperref[def:real-lower-bound]{Lower Bound}
\end{itemize}
\end{dependencies}

\begin{propositionbox}{Proposition (Positive Dilation Preserves Upper Bounds)}
\begin{proposition}[Positive Dilation Preserves Upper Bounds]
\label{prop:positive-dilation-preserves-upper-bounds}
Let $A\subseteq\mathbb{R}$, let $\lambda>0$, and let $u\in\mathbb{R}$.
If $u$ is an upper bound of $A$, then $\lambda u$ is an upper bound of
$\lambda A$; that is,
\[
\left(\forall a\in A,\ a\leq u\right)
\Longrightarrow
\left(\forall x\in\lambda A,\ x\leq \lambda u\right).
\]
\hyperref[prf:positive-dilation-preserves-upper-bounds]{Go to proof.}
\end{proposition}
\end{propositionbox}

\begin{remark*}[Standard quantified statement]
\[
\lambda>0\land(\forall a\in A)(a\leq u)
\Longrightarrow
(\forall x\in\lambda A)(x\leq\lambda u).
\]
\end{remark*}

\begin{remark*}[Predicate reading]
\[
\begin{aligned}
&P=\mathsf{OrderedSet}(\mathbb{R},\leq),\qquad
A\subseteq\mathbb{R},\qquad \lambda A=\{\lambda a:a\in A\},\\
&\lambda>0\land\operatorname{UpperBound}(u,A,P)
\Longrightarrow\operatorname{UpperBound}(\lambda u,\lambda A,P).
\end{aligned}
\]
\end{remark*}

\begin{remark*}[Negated quantified statement]
\[
\lambda>0,\qquad(\forall a\in A)(a\leq u),\qquad
(\exists y\in\lambda A)(\lambda u<y).
\]
\end{remark*}

\begin{remark*}[Negation predicate reading]
\[
\begin{aligned}
&P=\mathsf{OrderedSet}(\mathbb{R},\leq),\qquad
A\subseteq\mathbb{R},\qquad \lambda A=\{\lambda a:a\in A\},\\
&\lambda>0\land\operatorname{UpperBound}(u,A,P)\land
\neg\operatorname{UpperBound}(\lambda u,\lambda A,P).
\end{aligned}
\]
\end{remark*}

\begin{remark*}[Failure modes]
\begin{description}
\item[Exposition.]
A failure would require positive multiplication to reverse an inequality.


  \item[\operatorname{UpperBound}.]
  \[
\lambda>0,\qquad(\forall a\in A)(a\leq u),\qquad
(\exists y\in\lambda A)(\lambda u<y).
\]
\[
\begin{aligned}
&P=\mathsf{OrderedSet}(\mathbb{R},\leq),\qquad
A\subseteq\mathbb{R},\qquad \lambda A=\{\lambda a:a\in A\},\\
&\lambda>0\land\operatorname{UpperBound}(u,A,P)\land
\neg\operatorname{UpperBound}(\lambda u,\lambda A,P).
\end{aligned}
\]
\end{description}
\end{remark*}

\begin{remark*}[Failure modes]
\begin{description}
\item[Exposition.]
Write $y=\lambda a$ with $a\in A$.  Since $a\leq u$ and $\lambda>0$,
$\lambda a\leq\lambda u$.


  \item[\operatorname{UpperBound}.]
  \[
\lambda>0,\qquad(\forall a\in A)(a\leq u),\qquad
(\exists y\in\lambda A)(\lambda u<y).
\]
\[
\begin{aligned}
&P=\mathsf{OrderedSet}(\mathbb{R},\leq),\qquad
A\subseteq\mathbb{R},\qquad \lambda A=\{\lambda a:a\in A\},\\
&\lambda>0\land\operatorname{UpperBound}(u,A,P)\land
\neg\operatorname{UpperBound}(\lambda u,\lambda A,P).
\end{aligned}
\]
\end{description}
\end{remark*}

\begin{remark*}[Interpretation]
Positive dilation stretches or contracts the set without reversing order, so
upper bounds dilate to upper bounds.
\end{remark*}

\begin{dependencies}
\begin{itemize}
  \item \hyperref[def:reals]{Real Numbers}
  \item \hyperref[def:ordered-set]{Ordered Set}
  \item \hyperref[def:subset]{Subset}
  \item \hyperref[def:multiplication-on-r]{Multiplication on $\mathbb{R}$}
  \item \hyperref[def:displacement-dilation-reflection]{Displacement, Dilation, and Reflection}
  \item \hyperref[def:real-upper-bound]{Upper Bound}
\end{itemize}
\end{dependencies}

\begin{propositionbox}{Proposition (Positive Dilation Preserves Lower Bounds)}
\begin{proposition}[Positive Dilation Preserves Lower Bounds]
\label{prop:positive-dilation-preserves-lower-bounds}
Let $A\subseteq\mathbb{R}$, let $\lambda>0$, and let
$\ell\in\mathbb{R}$.  If $\ell$ is a lower bound of $A$, then
$\lambda\ell$ is a lower bound of $\lambda A$; that is,
\[
\left(\forall a\in A,\ \ell\leq a\right)
\Longrightarrow
\left(\forall x\in\lambda A,\ \lambda\ell\leq x\right).
\]
\hyperref[prf:positive-dilation-preserves-lower-bounds]{Go to proof.}
\end{proposition}
\end{propositionbox}

\begin{remark*}[Standard quantified statement]
\[
\lambda>0\land(\forall a\in A)(\ell\leq a)
\Longrightarrow
(\forall x\in\lambda A)(\lambda\ell\leq x).
\]
\end{remark*}

\begin{remark*}[Predicate reading]
\[
\begin{aligned}
&P=\mathsf{OrderedSet}(\mathbb{R},\leq),\qquad
A\subseteq\mathbb{R},\qquad \lambda A=\{\lambda a:a\in A\},\\
&\lambda>0\land\operatorname{LowerBound}(\ell,A,P)
\Longrightarrow\operatorname{LowerBound}(\lambda\ell,\lambda A,P).
\end{aligned}
\]
\end{remark*}

\begin{remark*}[Negated quantified statement]
\[
\lambda>0,\qquad(\forall a\in A)(\ell\leq a),\qquad
(\exists y\in\lambda A)(y<\lambda\ell).
\]
\end{remark*}

\begin{remark*}[Negation predicate reading]
\[
\begin{aligned}
&P=\mathsf{OrderedSet}(\mathbb{R},\leq),\qquad
A\subseteq\mathbb{R},\qquad \lambda A=\{\lambda a:a\in A\},\\
&\lambda>0\land\operatorname{LowerBound}(\ell,A,P)\land
\neg\operatorname{LowerBound}(\lambda\ell,\lambda A,P).
\end{aligned}
\]
\end{remark*}

\begin{remark*}[Failure modes]
\begin{description}
\item[Exposition.]
A failure would require a positively dilated point below the positively
dilated lower bound.


  \item[\operatorname{LowerBound}.]
  \[
\lambda>0,\qquad(\forall a\in A)(\ell\leq a),\qquad
(\exists y\in\lambda A)(y<\lambda\ell).
\]
\[
\begin{aligned}
&P=\mathsf{OrderedSet}(\mathbb{R},\leq),\qquad
A\subseteq\mathbb{R},\qquad \lambda A=\{\lambda a:a\in A\},\\
&\lambda>0\land\operatorname{LowerBound}(\ell,A,P)\land
\neg\operatorname{LowerBound}(\lambda\ell,\lambda A,P).
\end{aligned}
\]
\end{description}
\end{remark*}

\begin{remark*}[Failure modes]
\begin{description}
\item[Exposition.]
Write $y=\lambda a$ with $a\in A$.  Since $\ell\leq a$ and $\lambda>0$,
$\lambda\ell\leq\lambda a$.


  \item[\operatorname{LowerBound}.]
  \[
\lambda>0,\qquad(\forall a\in A)(\ell\leq a),\qquad
(\exists y\in\lambda A)(y<\lambda\ell).
\]
\[
\begin{aligned}
&P=\mathsf{OrderedSet}(\mathbb{R},\leq),\qquad
A\subseteq\mathbb{R},\qquad \lambda A=\{\lambda a:a\in A\},\\
&\lambda>0\land\operatorname{LowerBound}(\ell,A,P)\land
\neg\operatorname{LowerBound}(\lambda\ell,\lambda A,P).
\end{aligned}
\]
\end{description}
\end{remark*}

\begin{remark*}[Interpretation]
Positive dilation preserves lower-bound information for the same reason it
preserves upper-bound information: order is not reversed.
\end{remark*}

\begin{dependencies}
\begin{itemize}
  \item \hyperref[def:reals]{Real Numbers}
  \item \hyperref[def:ordered-set]{Ordered Set}
  \item \hyperref[def:subset]{Subset}
  \item \hyperref[def:multiplication-on-r]{Multiplication on $\mathbb{R}$}
  \item \hyperref[def:displacement-dilation-reflection]{Displacement, Dilation, and Reflection}
  \item \hyperref[def:real-lower-bound]{Lower Bound}
\end{itemize}
\end{dependencies}

\begin{propositionbox}{Proposition (Negative Dilation Sends Lower Bounds to Upper Bounds)}
\begin{proposition}[Negative Dilation Sends Lower Bounds to Upper Bounds]
\label{prop:negative-dilation-sends-lower-to-upper-bounds}
Let $A\subseteq\mathbb{R}$, let $\lambda<0$, and let
$\ell\in\mathbb{R}$.  If $\ell$ is a lower bound of $A$, then
$\lambda\ell$ is an upper bound of $\lambda A$; that is,
\[
\left(\forall a\in A,\ \ell\leq a\right)
\Longrightarrow
\left(\forall x\in\lambda A,\ x\leq \lambda\ell\right).
\]
\hyperref[prf:negative-dilation-sends-lower-to-upper-bounds]{Go to proof.}
\end{proposition}
\end{propositionbox}

\begin{remark*}[Standard quantified statement]
\[
\lambda<0\land(\forall a\in A)(\ell\leq a)
\Longrightarrow
(\forall x\in\lambda A)(x\leq\lambda\ell).
\]
\end{remark*}

\begin{remark*}[Predicate reading]
\[
\begin{aligned}
&P=\mathsf{OrderedSet}(\mathbb{R},\leq),\qquad
A\subseteq\mathbb{R},\qquad \lambda A=\{\lambda a:a\in A\},\\
&\lambda<0\land\operatorname{LowerBound}(\ell,A,P)
\Longrightarrow\operatorname{UpperBound}(\lambda\ell,\lambda A,P).
\end{aligned}
\]
\end{remark*}

\begin{remark*}[Negated quantified statement]
\[
\lambda<0,\qquad(\forall a\in A)(\ell\leq a),\qquad
(\exists y\in\lambda A)(\lambda\ell<y).
\]
\end{remark*}

\begin{remark*}[Negation predicate reading]
\[
\begin{aligned}
&P=\mathsf{OrderedSet}(\mathbb{R},\leq),\qquad
A\subseteq\mathbb{R},\qquad \lambda A=\{\lambda a:a\in A\},\\
&\lambda<0\land\operatorname{LowerBound}(\ell,A,P)\land
\neg\operatorname{UpperBound}(\lambda\ell,\lambda A,P).
\end{aligned}
\]
\end{remark*}

\begin{remark*}[Failure modes]
\begin{description}
\item[Exposition.]
A failure would require negative multiplication not to reverse the lower-bound
inequality.


  \item[\operatorname{LowerBound} / \operatorname{UpperBound}.]
  \[
\lambda<0,\qquad(\forall a\in A)(\ell\leq a),\qquad
(\exists y\in\lambda A)(\lambda\ell<y).
\]
\[
\begin{aligned}
&P=\mathsf{OrderedSet}(\mathbb{R},\leq),\qquad
A\subseteq\mathbb{R},\qquad \lambda A=\{\lambda a:a\in A\},\\
&\lambda<0\land\operatorname{LowerBound}(\ell,A,P)\land
\neg\operatorname{UpperBound}(\lambda\ell,\lambda A,P).
\end{aligned}
\]
\end{description}
\end{remark*}

\begin{remark*}[Failure modes]
\begin{description}
\item[Exposition.]
Write $y=\lambda a$ with $a\in A$.  Since $\ell\leq a$ and $\lambda<0$,
$\lambda a\leq\lambda\ell$.


  \item[\operatorname{LowerBound} / \operatorname{UpperBound}.]
  \[
\lambda<0,\qquad(\forall a\in A)(\ell\leq a),\qquad
(\exists y\in\lambda A)(\lambda\ell<y).
\]
\[
\begin{aligned}
&P=\mathsf{OrderedSet}(\mathbb{R},\leq),\qquad
A\subseteq\mathbb{R},\qquad \lambda A=\{\lambda a:a\in A\},\\
&\lambda<0\land\operatorname{LowerBound}(\ell,A,P)\land
\neg\operatorname{UpperBound}(\lambda\ell,\lambda A,P).
\end{aligned}
\]
\end{description}
\end{remark*}

\begin{remark*}[Interpretation]
Negative dilation reverses the order on the real line, so lower bounds become
upper bounds after dilation.
\end{remark*}

\begin{dependencies}
\begin{itemize}
  \item \hyperref[def:reals]{Real Numbers}
  \item \hyperref[def:ordered-set]{Ordered Set}
  \item \hyperref[def:subset]{Subset}
  \item \hyperref[def:multiplication-on-r]{Multiplication on $\mathbb{R}$}
  \item \hyperref[def:displacement-dilation-reflection]{Displacement, Dilation, and Reflection}
  \item \hyperref[def:real-lower-bound]{Lower Bound}
  \item \hyperref[def:real-upper-bound]{Upper Bound}
\end{itemize}
\end{dependencies}

\begin{propositionbox}{Proposition (Negative Dilation Sends Upper Bounds to Lower Bounds)}
\begin{proposition}[Negative Dilation Sends Upper Bounds to Lower Bounds]
\label{prop:negative-dilation-sends-upper-to-lower-bounds}
Let $A\subseteq\mathbb{R}$, let $\lambda<0$, and let $u\in\mathbb{R}$.
If $u$ is an upper bound of $A$, then $\lambda u$ is a lower bound of
$\lambda A$; that is,
\[
\left(\forall a\in A,\ a\leq u\right)
\Longrightarrow
\left(\forall x\in\lambda A,\ \lambda u\leq x\right).
\]
\hyperref[prf:negative-dilation-sends-upper-to-lower-bounds]{Go to proof.}
\end{proposition}
\end{propositionbox}

\begin{remark*}[Standard quantified statement]
\[
\lambda<0\land(\forall a\in A)(a\leq u)
\Longrightarrow
(\forall x\in\lambda A)(\lambda u\leq x).
\]
\end{remark*}

\begin{remark*}[Predicate reading]
\[
\begin{aligned}
&P=\mathsf{OrderedSet}(\mathbb{R},\leq),\qquad
A\subseteq\mathbb{R},\qquad \lambda A=\{\lambda a:a\in A\},\\
&\lambda<0\land\operatorname{UpperBound}(u,A,P)
\Longrightarrow\operatorname{LowerBound}(\lambda u,\lambda A,P).
\end{aligned}
\]
\end{remark*}

\begin{remark*}[Negated quantified statement]
\[
\lambda<0,\qquad(\forall a\in A)(a\leq u),\qquad
(\exists y\in\lambda A)(y<\lambda u).
\]
\end{remark*}

\begin{remark*}[Negation predicate reading]
\[
\begin{aligned}
&P=\mathsf{OrderedSet}(\mathbb{R},\leq),\qquad
A\subseteq\mathbb{R},\qquad \lambda A=\{\lambda a:a\in A\},\\
&\lambda<0\land\operatorname{UpperBound}(u,A,P)\land
\neg\operatorname{LowerBound}(\lambda u,\lambda A,P).
\end{aligned}
\]
\end{remark*}

\begin{remark*}[Failure modes]
\begin{description}
\item[Exposition.]
A failure would require negative multiplication not to reverse the upper-bound
inequality.


  \item[\operatorname{UpperBound} / \operatorname{LowerBound}.]
  \[
\lambda<0,\qquad(\forall a\in A)(a\leq u),\qquad
(\exists y\in\lambda A)(y<\lambda u).
\]
\[
\begin{aligned}
&P=\mathsf{OrderedSet}(\mathbb{R},\leq),\qquad
A\subseteq\mathbb{R},\qquad \lambda A=\{\lambda a:a\in A\},\\
&\lambda<0\land\operatorname{UpperBound}(u,A,P)\land
\neg\operatorname{LowerBound}(\lambda u,\lambda A,P).
\end{aligned}
\]
\end{description}
\end{remark*}

\begin{remark*}[Failure modes]
\begin{description}
\item[Exposition.]
Write $y=\lambda a$ with $a\in A$.  Since $a\leq u$ and $\lambda<0$,
$\lambda u\leq\lambda a$.


  \item[\operatorname{UpperBound} / \operatorname{LowerBound}.]
  \[
\lambda<0,\qquad(\forall a\in A)(a\leq u),\qquad
(\exists y\in\lambda A)(y<\lambda u).
\]
\[
\begin{aligned}
&P=\mathsf{OrderedSet}(\mathbb{R},\leq),\qquad
A\subseteq\mathbb{R},\qquad \lambda A=\{\lambda a:a\in A\},\\
&\lambda<0\land\operatorname{UpperBound}(u,A,P)\land
\neg\operatorname{LowerBound}(\lambda u,\lambda A,P).
\end{aligned}
\]
\end{description}
\end{remark*}

\begin{remark*}[Interpretation]
Negative dilation sends upper-bound information to lower-bound information
because multiplying by a negative scalar reverses inequalities.
\end{remark*}

\begin{dependencies}
\begin{itemize}
  \item \hyperref[def:reals]{Real Numbers}
  \item \hyperref[def:ordered-set]{Ordered Set}
  \item \hyperref[def:subset]{Subset}
  \item \hyperref[def:multiplication-on-r]{Multiplication on $\mathbb{R}$}
  \item \hyperref[def:displacement-dilation-reflection]{Displacement, Dilation, and Reflection}
  \item \hyperref[def:real-upper-bound]{Upper Bound}
  \item \hyperref[def:real-lower-bound]{Lower Bound}
\end{itemize}
\end{dependencies}

\begin{corollarybox}{Corollary (Reflection Swaps Upper and Lower Bounds)}
\begin{corollary}[Reflection Swaps Upper and Lower Bounds]
\label{cor:reflection-swaps-upper-lower-bounds}
Let $A\subseteq\mathbb{R}$. For every $b\in\mathbb{R}$,
\[
b\text{ is an upper bound of }A
\Longleftrightarrow
-b\text{ is a lower bound of }-A,
\]
and
\[
b\text{ is a lower bound of }A
\Longleftrightarrow
-b\text{ is an upper bound of }-A.
\]
Equivalently,
\[
\left(\forall a\in A,\ a\leq b\right)
\Longleftrightarrow
\left(\forall x\in -A,\ -b\leq x\right),
\]
and
\[
\left(\forall a\in A,\ b\leq a\right)
\Longleftrightarrow
\left(\forall x\in -A,\ x\leq -b\right).
\]
\hyperref[prf:reflection-swaps-upper-lower-bounds]{Go to proof.}
\end{corollary}
\end{corollarybox}

\begin{remark*}[Standard quantified statement]
\[
(\forall a\in A)(a\leq b)\Longleftrightarrow
(\forall x\in -A)(-b\leq x),
\]
and
\[
(\forall a\in A)(b\leq a)\Longleftrightarrow
(\forall x\in -A)(x\leq -b).
\]
\end{remark*}

\begin{remark*}[Predicate reading]
\[
\begin{aligned}
&P=\mathsf{OrderedSet}(\mathbb{R},\leq),\qquad
A\subseteq\mathbb{R},\qquad -A=\{-a:a\in A\},\\
&\operatorname{UpperBound}(b,A,P)\Longleftrightarrow\operatorname{LowerBound}(-b,-A,P),\\
&\operatorname{LowerBound}(b,A,P)\Longleftrightarrow\operatorname{UpperBound}(-b,-A,P).
\end{aligned}
\]
\end{remark*}

\begin{remark*}[Negated quantified statement]
A failure is a failure of one of the displayed biconditionals:
\[
\bigl[(\forall a\in A)(a\leq b)\bigr]
\not\Longleftrightarrow
\bigl[(\forall x\in -A)(-b\leq x)\bigr],
\]
or
\[
\bigl[(\forall a\in A)(b\leq a)\bigr]
\not\Longleftrightarrow
\bigl[(\forall x\in -A)(x\leq -b)\bigr].
\]
\end{remark*}

\begin{remark*}[Negation predicate reading]
\[
\begin{aligned}
&P=\mathsf{OrderedSet}(\mathbb{R},\leq),\qquad
A\subseteq\mathbb{R},\qquad -A=\{-a:a\in A\},\\
&\operatorname{UpperBound}(b,A,P)\not\Longleftrightarrow\operatorname{LowerBound}(-b,-A,P)
\end{aligned}
\]
or
\[
\begin{aligned}
&P=\mathsf{OrderedSet}(\mathbb{R},\leq),\qquad
A\subseteq\mathbb{R},\qquad -A=\{-a:a\in A\},\\
&\operatorname{LowerBound}(b,A,P)\not\Longleftrightarrow\operatorname{UpperBound}(-b,-A,P).
\end{aligned}
\]
\end{remark*}

\begin{remark*}[Failure modes]
\begin{description}
\item[Exposition.]
Reflection reverses the real line, so lower information becomes upper
information and upper information becomes lower information.


  \item[\operatorname{LowerBound} / \operatorname{UpperBound}.]

\[
\begin{aligned}
&P=\mathsf{OrderedSet}(\mathbb{R},\leq),\qquad
A\subseteq\mathbb{R},\qquad -A=\{-a:a\in A\},\\
&\operatorname{UpperBound}(b,A,P)\not\Longleftrightarrow\operatorname{LowerBound}(-b,-A,P)
\end{aligned}
\]
\[
\begin{aligned}
&P=\mathsf{OrderedSet}(\mathbb{R},\leq),\qquad
A\subseteq\mathbb{R},\qquad -A=\{-a:a\in A\},\\
&\operatorname{UpperBound}(b,A,P)\Longleftrightarrow\operatorname{LowerBound}(-b,-A,P),\\
&\operatorname{LowerBound}(b,A,P)\Longleftrightarrow\operatorname{UpperBound}(-b,-A,P).
\end{aligned}
\]
\end{description}
\end{remark*}

\begin{remark*}[Failure modes]
\begin{description}
\item[Exposition.]
Apply negative dilation with $\lambda=-1$.


  \item[\operatorname{LowerBound} / \operatorname{UpperBound}.]

\[
\begin{aligned}
&P=\mathsf{OrderedSet}(\mathbb{R},\leq),\qquad
A\subseteq\mathbb{R},\qquad -A=\{-a:a\in A\},\\
&\operatorname{LowerBound}(b,A,P)\not\Longleftrightarrow\operatorname{UpperBound}(-b,-A,P)
\end{aligned}
\]
\[
\begin{aligned}
&P=\mathsf{OrderedSet}(\mathbb{R},\leq),\qquad
A\subseteq\mathbb{R},\qquad -A=\{-a:a\in A\},\\
&\operatorname{UpperBound}(b,A,P)\Longleftrightarrow\operatorname{LowerBound}(-b,-A,P),\\
&\operatorname{LowerBound}(b,A,P)\Longleftrightarrow\operatorname{UpperBound}(-b,-A,P).
\end{aligned}
\]
\end{description}
\end{remark*}

\begin{remark*}[Interpretation]
Reflection is negative dilation by $-1$, so it exchanges the roles of upper
and lower bounds.
\end{remark*}

\begin{dependencies}
\begin{itemize}
  \item \hyperref[def:reals]{Real Numbers}
  \item \hyperref[def:ordered-set]{Ordered Set}
  \item \hyperref[def:subset]{Subset}
  \item \hyperref[def:multiplication-on-r]{Multiplication on $\mathbb{R}$}
  \item \hyperref[def:displacement-dilation-reflection]{Displacement, Dilation, and Reflection}
  \item \hyperref[def:real-lower-bound]{Lower Bound}
  \item \hyperref[def:real-upper-bound]{Upper Bound}
  \item \hyperref[prop:negative-dilation-sends-lower-to-upper-bounds]{Negative Dilation Sends Lower Bounds to Upper Bounds}
  \item \hyperref[prop:negative-dilation-sends-upper-to-lower-bounds]{Negative Dilation Sends Upper Bounds to Lower Bounds}
\end{itemize}
\end{dependencies}
